\chapter{Reproducibility}
\label{app:reproducibility}

\section{Report Build and Asset Identity}

The report source includes its bibliography and selected figures and tables. With the required LaTeX tools installed, these
files support a PDF build without rerunning an experiment or contacting a model provider. The assets are bundled with the
separately retained report source. The experiment artefacts remain in the development workspace and are not needed for this build.

The asset register records each bundled file's purpose, hash and review status. The Python publication code creates candidate
outputs in the development workspace. After review, the selected PDF or table is promoted into the relevant chapter asset
folder and its register entry is updated. This records file identity. It does not establish scientific correctness,
redistribution permission or final visual approval.

The acquisition exports with superseded labels remain outside the report. The results chapter contains only the tables and
figures selected for the current draft. The system-boundary, study-sequence and CSEDM entries now point to corrected exports
generated from committed source revisions.

Build the report from the \texttt{report/} directory:

\begin{verbatim}
latexmk -pdf main.tex
\end{verbatim}

This command invokes \texttt{latexmk} with PDF output, non-interactive error reporting and automated multi-pass
dependency resolution across chapters, Biber, and vector graphics. The required tools are \texttt{latexmk}, Biber and
the installed TeX packages. The report README records the compilation commands.

A retained report archive labelled \texttt{1639b14} was built separately on 11 September 2026 using the installed TeX tools.
It produced 63 pages (49 content pages), and its recorded inputs contained no path into the working directory or experiment
artefacts. The archive, PDF and build records were retained locally with their checksums. The current repository no longer
resolves that identifier to a report tree, so this is recovery evidence for the retained archive rather than independently
verifiable Git provenance. It does not cover the current report source or a fresh installation of the required tools.

\section{Analysis Reproduction and Restricted Inputs}

The implementation uses a locked Pixi environment. Retained model responses, execution records, analysis plans and reports
allow offline inspection without purchasing another generation. Replaying recorded responses checks processing of those
responses. It does not establish that a fresh provider request will produce identical text. Re-executing generated code also
has different prerequisites from reading a saved sandbox result.

The development source retained with that earlier checkpoint was also restored separately. Three existing API tests passed
using the restored source and installed dependencies, covering both practice tasks and the recorded prediction-and-reveal flow.
The offline development benchmark command completed with eight condition predictions, two criterion records and five
Parquet datasets. Its outcomes are authored fixtures. These checks establish that the bundled development inputs support
those workflows. Browser behaviour and installation into a fresh environment require further verification.

The acquisition primary analysis was reproduced on 11 September 2026 in a separate local clone at revision
\texttt{5d6e76b}. Its historical environment was installed with \texttt{pixi install --locked}, using available package
caches. The analysis read the retained 98,000 public records. All six generated files matched the original files byte for
byte, including the report and comparison tables. No model or sandbox calls were made. This verifies the primary analysis
from recorded observations. Regeneration of the simulation and installation on another machine remain unverified.

The tracker simulation and analysis were also regenerated in a separate locked environment at surviving revision
\texttt{d14d8a2}. The trajectories matched the retained file byte for byte. Every reported metric, interval and decision
agreed. The revision fields and dependent hashes differed. The original recorded revision, \texttt{baf40df}, could not be
recovered from the current repository or its available rewrite map. This check establishes numerical reproduction from
the named surviving implementation, while the link to the original source revision remains unresolved.

The CSEDM pipeline requires authorised access to the local archive identified in its inventory configuration. Its official
documentation describes the DataShop access route, prediction cut-off and supplied learner-separated folds
\parencite[Accessing the Data, Evaluation and The Data]{csedm2019challenge}. Download access must not be treated as permission to
redistribute learner records. The public package excludes raw events, learner identifiers and student source code. Rebuilding
the analysis therefore requires inputs beyond a public checkout.

From a clean committed development checkout, with the locked environment installed and the configured archive available,
the companion pipeline is:

\begin{verbatim}
cd development
pixi run python scripts/csedm_study_reproduce.py
\end{verbatim}

The command rebuilds inventory, features, predictions, analysis and publication, then reports the output locations. Its default
output directory includes the code revision. A reproduction attempt that changes output content requires a separately named
directory. Supplying an arbitrary revision string
does not establish that the running code matches that revision.

Compare the resulting scientific quantities, sample counts and input hashes with the retained reference reports. A new code
revision can change provenance fields and manifest hashes even when numerical results agree. Record those differences rather
than promising that every new build reproduces an old manifest byte for byte.

\section{Scope of the Verification}

The recorded checks cover isolated PDF compilation, offline benchmark fixtures, and numerical replay of the retained acquisition
and tracker analyses. They do not cover installation without the available package caches, installation on another machine, new provider
generation, or redistribution of restricted learner records. A restore on the same machine supplies useful recovery evidence.
The examiner package must also retain the final report archive and checksum, release revision, lockfile identity, reference
reports, unresolved access requirements, and a verified off-device backup.
